\setlength{\absparsep}{18pt}
\begin{resumo}
Em janeiro de 2019 o Ceará tornou-se o primeiro estado brasileiro a
proibir a pulverização aérea de agrotóxicos. A lei respondeu a um conflito que vinha dos
anos 2000 na Chapada do Apodi, onde bananais de exportação eram pulverizados por
avião ao lado de comunidades rurais. Este projeto propõe avaliar o impacto da
proibição sobre a saúde ao nascer. O desenho é de diferenças em diferenças na
fronteira: municípios cearenses da Chapada, atingidos pela lei, são comparados
com municípios potiguares da mesma chapada, que partilham clima e aquífero e não foram atingidos. O estimador é o de Callaway e Sant'Anna, com uma
coorte tratada em 2019, e o desfecho primário é o peso ao nascer registrado no
SINASC. O projeto organiza a avaliação pela teoria da mudança da política,
explicita o grafo causal e as hipóteses de identificação, e mostra, com o
cálculo de poder, o limite do que se pode aprender: o efeito esperado na média
municipal é pequeno, porque a deriva alcança poucas centenas de metros. Por
isso o plano acrescenta ao contraste principal um desenho de calendário,
desfechos de perda fetal e a medição, por pedidos de acesso à informação,
de quanto avião havia antes da lei. Um segundo ensaio usa os resultados para um
custo-benefício da proibição e para discutir a escolha do instrumento.

\vspace{\onelineskip}
\noindent\textbf{Palavras-chave}: agrotóxicos; pulverização aérea; avaliação de
impacto; diferenças em diferenças; saúde ao nascer; Chapada do Apodi.
\end{resumo}

\begin{resumo}[Abstract]
\begin{otherlanguage*}{english}
In January 2019 Ceará became the first Brazilian state to ban the
aerial spraying of pesticides. The law followed a conflict dating from the 2000s in the
Chapada do Apodi, where export banana plantations were sprayed by aircraft next
to rural communities. This project proposes an impact evaluation of the ban on
health at birth. The design is a border difference-in-differences: municipalities
on the Ceará side of the plateau, covered by the law, are compared with
municipalities on the Rio Grande do Norte side, which share climate and aquifer and were not covered. The estimator is Callaway and Sant'Anna's, with a
single cohort treated in 2019, and the primary outcome is birth weight from the
national birth registry. The project organises the evaluation around the
policy's theory of change, states the causal graph and the identifying
assumptions, and uses a power calculation to show the limits of what can be
learned: the expected effect on the municipal average is small, because drift
reaches only a few hundred metres. The plan therefore adds to the main contrast
a calendar design, fetal-loss outcomes, and freedom-of-information requests that
measure how much aerial spraying existed before the law. A second
essay uses the results for a cost-benefit analysis of the ban and a discussion
of instrument choice.

\vspace{\onelineskip}
\noindent\textbf{Keywords}: pesticides; aerial spraying; impact evaluation;
difference-in-differences; health at birth; Chapada do Apodi.
\end{otherlanguage*}
\end{resumo}
